\documentclass[preview]{standalone}

\usepackage{amsmath}
\usepackage{amssymb}
\usepackage{stellar}
\usepackage{definitions}

\begin{document}

\id{minimal-polynomial-canonical-forms}
\genpage

\section{Cyclic and block-diagonal matrices}

\begin{snippetproposition}{cyclic-endomorphism-minimal-polynomial}{Minimal polynomial of a cyclic endomorphism}
    Let \(V=K[X]v\) be the cyclic module associated with a
    \linearendomorphism \(\varphi\), and suppose
    \[
        \moduleann(v)=(p)
    \]
    for a monic polynomial \(p\). Then every matrix \(A\) representing
    \(\varphi\) satisfies
    \[
        \mu_A=\chi_A=p.
    \]
\end{snippetproposition}

\begin{snippetproof}{cyclic-endomorphism-minimal-polynomial-proof}{cyclic-endomorphism-minimal-polynomial}{Minimal polynomial of a cyclic endomorphism}
    In a cyclic basis, the matrix of \(\varphi\) is \(C(p)\), whose
    characteristic polynomial is \(p\). Moreover, a polynomial \(f\)
    annihilates \(\varphi\) only if it annihilates the cyclic generator
    \(v\), hence only if \(p\divides f\). Since \(p\) annihilates the
    whole quotient module \(K[X]/(p)\), it is the minimal polynomial.
    Similarity invariance extends the conclusion to every representing
    matrix.
\end{snippetproof}

\begin{snippetcorollary}{companion-matrix-minimal-polynomial}{Minimal polynomial of a companion matrix}
    If \(p\in K[X]\) is monic, then
    \[
        \mu_{C(p)}=\chi_{C(p)}=p.
    \]
\end{snippetcorollary}

\begin{snippetproof}{companion-matrix-minimal-polynomial-proof}{companion-matrix-minimal-polynomial}{Minimal polynomial of a companion matrix}
    Apply the cyclic-endomorphism result to
    \(K[X]/(p)\), with generator \(1+(p)\).
\end{snippetproof}

\begin{snippetcorollary}{jordan-block-minimal-polynomial}{Minimal polynomial of a Jordan block}
    For every \jordanblock \(J_n(a)\),
    \[
        \mu_{J_n(a)}=\chi_{J_n(a)}=(X-a)^n.
    \]
\end{snippetcorollary}

\begin{snippetproof}{jordan-block-minimal-polynomial-proof}{jordan-block-minimal-polynomial}{Minimal polynomial of a Jordan block}
    The block \(J_n(a)\) represents the action of \(X\) on the cyclic
    module \(K[X]/((X-a)^n)\).
\end{snippetproof}

\begin{snippetproposition}{block-diagonal-minimal-polynomial-lcm}{Minimal polynomial of a block-diagonal matrix}
    Let
    \[
        B=\operatorname{diag}(B_1,\ldots,B_s).
    \]
    For every \(f\in K[X]\),
    \[
        f(B)=\operatorname{diag}(f(B_1),\ldots,f(B_s)).
    \]
    Consequently,
    \[
        \mu_B=\lcm(\mu_{B_1},\ldots,\mu_{B_s}),
    \]
    where the least common multiple is chosen monic.
\end{snippetproposition}

\begin{snippetproof}{block-diagonal-minimal-polynomial-lcm-proof}{block-diagonal-minimal-polynomial-lcm}{Minimal polynomial of a block-diagonal matrix}
    Every power of a block-diagonal matrix is obtained by taking the
    corresponding power of each block, which proves the formula for
    \(f(B)\). Thus \(f(B)=0\) exactly when every \(\mu_{B_i}\) divides
    \(f\). The monic polynomial of least degree with this property is
    their least common multiple.
\end{snippetproof}

\section{Canonical forms}

\begin{snippettheorem}{minimal-polynomial-rational-canonical-form}{Minimal polynomial from rational canonical form}
    Suppose the \rationalcanonicalform of \(A\) has companion blocks
    \[
        C(d_1),\ldots,C(d_s),
        \qquad
        d_1\divides\cdots\divides d_s.
    \]
    Then
    \[
        \mu_A=d_s,
        \qquad
        \chi_A=d_1\cdots d_s.
    \]
\end{snippettheorem}

\begin{snippetproof}{minimal-polynomial-rational-canonical-form-proof}{minimal-polynomial-rational-canonical-form}{Minimal polynomial from rational canonical form}
    The minimal polynomial is the least common multiple of the minimal
    polynomials \(d_i\) of the companion blocks. The divisibility chain
    makes this least common multiple equal to \(d_s\). The
    characteristic-polynomial formula is the product formula for
    companion blocks.
\end{snippetproof}

\begin{snippetcorollary}{minimal-characteristic-polynomial-irreducible-factors}{Minimal and characteristic polynomials have the same irreducible factors}
    The \minimalpolynomial of \(A\) divides its
    \characteristicpolynomial, and the two polynomials have exactly the
    same monic irreducible factors. In particular, one splits into
    linear factors over \(K\) \ifandonlyif the other does.
\end{snippetcorollary}

\begin{snippetproof}{minimal-characteristic-polynomial-irreducible-factors-proof}{minimal-characteristic-polynomial-irreducible-factors}{Minimal and characteristic polynomials have the same irreducible factors}
    In rational canonical form,
    \[
        \mu_A=d_s,
        \qquad
        \chi_A=d_1\cdots d_s,
        \qquad
        d_i\divides d_s.
    \]
    Thus \(\mu_A\divides\chi_A\). Every irreducible factor of \(\mu_A\)
    appears in \(\chi_A\), while every irreducible factor of some
    \(d_i\), and hence of \(\chi_A\), divides \(d_s=\mu_A\).
\end{snippetproof}

\begin{snippettheorem}{minimal-polynomial-jordan-form}{Minimal polynomial from Jordan form}
    Suppose the characteristic polynomial of \(A\) splits over \(K\).
    For each distinct \eigenvalue \(a_i\), let \(n_i\) be the maximum
    size of a Jordan block belonging to \(a_i\). Then
    \[
        \mu_A
        =
        \prod_i(X-a_i)^{n_i}.
    \]
\end{snippettheorem}

\begin{snippetproof}{minimal-polynomial-jordan-form-proof}{minimal-polynomial-jordan-form}{Minimal polynomial from Jordan form}
    By similarity invariance, compute on the Jordan form. The minimal
    polynomial of \(J_m(a)\) is \((X-a)^m\), and the minimal polynomial
    of a block-diagonal matrix is the least common multiple of the
    block minimal polynomials. For a fixed eigenvalue this selects the
    largest exponent; factors belonging to distinct eigenvalues are
    coprime, so their least common multiple is the displayed product.
\end{snippetproof}

\section{Cayley--Hamilton and diagonalizability}

\begin{snippettheorem}{cayley-hamilton-theorem}{Cayley--Hamilton theorem}
    Every square \matrix over a \field annihilates its own
    \characteristicpolynomial:
    \[
        \chi_A(A)=0.
    \]
\end{snippettheorem}

\begin{snippetproof}{cayley-hamilton-theorem-proof}{cayley-hamilton-theorem}{Cayley--Hamilton theorem}
    Rational canonical form shows that
    \(\mu_A\divides\chi_A\). Since every multiple of the minimal
    polynomial annihilates \(A\), the characteristic polynomial
    annihilates \(A\).
\end{snippetproof}

\begin{snippetcorollary}{minimal-polynomial-diagonalizability-criterion}{Diagonalizability criterion from the minimal polynomial}
    A square \matrix \(A\) over \(K\) is \diagonalizable over \(K\)
    \ifandonlyif its \minimalpolynomial is a product of distinct monic
    linear factors in \(K[X]\).
\end{snippetcorollary}

\begin{snippetproof}{minimal-polynomial-diagonalizability-criterion-proof}{minimal-polynomial-diagonalizability-criterion}{Diagonalizability criterion from the minimal polynomial}
    If \(A\) is diagonalizable with distinct eigenvalues
    \(a_1,\ldots,a_s\), then
    \[
        \mu_A=\prod_{i=1}^{s}(X-a_i).
    \]
    Conversely, if the minimal polynomial is a product of distinct
    linear factors, the characteristic polynomial also splits. In the
    Jordan formula for \(\mu_A\), every maximal block size is therefore
    \(1\), so every Jordan block has size \(1\) and \(A\) is
    diagonalizable.
\end{snippetproof}

\section{Examples}

\begin{snippetexample}{minimal-polynomial-jordan-blocks-example}{Reading the minimal polynomial from Jordan blocks}
    If
    \[
        A\sim J_2(2)\oplus J_1(2)\oplus J_2(3),
    \]
    then the largest block for each eigenvalue has size \(2\). Hence
    \[
        \mu_A=(X-2)^2(X-3)^2,
        \qquad
        \chi_A=(X-2)^3(X-3)^2.
    \]
\end{snippetexample}

\begin{snippetexample}{minimal-characteristic-polynomials-different-example}{Different minimal and characteristic polynomials}
    If
    \[
        A\sim J_3(0)\oplus J_2(3)\oplus J_1(3),
    \]
    then
    \[
        \chi_A=X^3(X-3)^3,
        \qquad
        \mu_A=X^3(X-3)^2.
    \]
    The exponent of each linear factor in \(\mu_A\) is the largest
    Jordan-block size for the corresponding eigenvalue, whereas its
    exponent in \(\chi_A\) is the sum of all those block sizes.
\end{snippetexample}

\end{document}
